\ch{Chapter 3: Research Methodology}\label{chap:Chapter_3}



This chapter presents the research methodology used to investigate professional awareness and perceptions of Digital Engineering capabilities within the Information Technology (IT) and Information Assurance (IA) domains. The study uses a quantitative survey approach to collect data from IT and IA professionals, enabling systematic analysis of awareness levels, perceived value, and anticipated benefits of Digital Engineering practices for IA and IT service delivery. The methodology follows a systems engineering lifecycle approach, applying structured engineering principles to research design while supporting rigorous traceability between research questions, survey instruments, and analytical approaches.

\section{Research Design Overview}

This research uses a quantitative, cross-sectional survey design to address the three research questions established in Chapter 1. Survey methodology best fits the research questions for three reasons. First, the research questions focus on measuring professional awareness and perceptions across a broad population of IT and IA practitioners, requiring data collection from enough respondents to establish representative findings. Second, survey methodology enables standardized data collection that supports statistical analysis and generalization of results to the broader professional population. Third, the anonymous nature of survey research encourages candid responses about professional knowledge gaps and organizational capabilities without concerns about professional reputation or organizational disclosure.

The cross-sectional design captures professional perceptions at a single point in time, providing a snapshot of current awareness and perceived value of Digital Engineering capabilities within the IT and Information Assurance professional community. While this design does not enable longitudinal analysis of changing perceptions, it establishes baseline data that will inform future research directions and support strategic decision-making about Digital Engineering adoption initiatives.

\subsection{Justification for Survey Methodology}

Selecting survey methodology over alternative approaches reflects deliberate consideration of what this research seeks to accomplish and why that objective warrants investigation. Three questions merit explicit answers: why investigate perceptions rather than implementations, why survey methodology rather than case study or pilot implementation, and why perceived value warrants dissertation-level research.

\subsubsection{Why Perceptions Matter for Technology Adoption}

Technology adoption research consistently shows that perceived value influences adoption decisions regardless of demonstrated actual value. Both the Technology Acceptance Model and its extensions establish that perceived usefulness and perceived ease of use predict behavioral intention to adopt technologies. Professionals who do not perceive value in a capability will not advocate for its adoption within their organizations, regardless of technical merit demonstrated in other contexts. As the literature review in Chapter 2 documents, Digital Engineering has demonstrated value in aerospace and defense contexts. However, demonstrated value in one domain does not automatically transfer to adoption in another. IT and Information Assurance professionals operate in different organizational contexts, face different constraints, and hold different professional identities than aerospace systems engineers. Whether these professionals recognize potential value in Digital Engineering capabilities for their work is an open question that must be answered before implementation research becomes meaningful.

Furthermore, if professionals are unaware of Digital Engineering capabilities, no amount of demonstrated value will drive adoption because the capabilities will not enter consideration during tool selection, process design, or strategic planning. Establishing current awareness levels identifies whether education and communication initiatives are necessary precursors to adoption efforts. If professionals are aware but do not perceive value, the theoretical premise that Digital Engineering addresses recognized needs within these domains requires reconsideration before investing in implementation research.

\subsubsection{Why Not a Case Study or Implementation Pilot}

This research considered and rejected alternative approaches for specific reasons. A case study examining Digital Engineering implementation within a specific organization would provide rich contextual data about practical implementation challenges and outcomes. However, case study findings would reflect the particular organizational context, culture, technical environment, and implementation approach of that organization. The findings would not establish whether the broader professional community recognizes value in Digital Engineering or is aware of these capabilities. A successful case study might show technical feasibility without indicating whether adoption would occur beyond the studied organization.

An implementation pilot would face similar limitations while requiring access to organizational resources and authority to implement Digital Engineering tools and practices. Such a pilot would measure actual rather than perceived value, but would do so within a single organizational context that may not represent the broader population. Implementation research also presumes that the target professional community recognizes sufficient potential value to warrant the investment required for implementation. This presumption is precisely what the current research tests.

Survey research enables assessment across a broad population of practitioners, establishing baseline awareness and perceived value data that informs whether subsequent case study or implementation research would find receptive audiences. Survey methodology is the appropriate starting point for investigating a nascent application domain where professional awareness and perceptions remain unknown.

\subsubsection{Why Perceived Value Warrants Dissertation Research}

Whether perceived value merits dissertation-level investigation reflects a particular view of research contribution. By addressing a documented gap, this research contributes to knowledge: the academic literature contains no investigation of whether IT and Information Assurance professionals perceive value in Digital Engineering capabilities. As Chapter 2 establishes, the theoretical framework posits that Digital Engineering offers solutions to documented problems in these domains. Testing whether practitioners recognize this potential value provides empirical grounding for the theoretical framework.

If research reveals that professionals recognize value in Digital Engineering capabilities, this finding validates the theoretical premise and establishes foundation for subsequent implementation research. If research reveals limited awareness or skepticism, this finding challenges the theoretical framework and identifies barriers that must be addressed before adoption can occur. Either outcome advances understanding and provides actionable direction for researchers and practitioners. The study's contribution is empirical: it establishes evidence regarding professional perceptions that informs subsequent research and practice decisions, advancing that evidence in place of advocacy.

\subsection{Systems Engineering Approach to Research Design}

To show rigor in the design, execution, analysis, and documentation of results, this research follows a systems engineering lifecycle. Applying systems engineering principles to the dissertation project lends structure and credibility to the dissertation process through disciplined practice and supporting artifacts. A model captures and tracks the principal method of the dissertation effort, preserving traceability to decisions, citations, and artifacts. This is the same traceability that Digital Engineering brings to complex system development.

This choice is methodological rather than stylistic. A recurring threat to survey validity is that the link between the instrument and the research objectives rests on the author's assertion rather than on documented evidence. Instead, this study represents each research question as a \emph{requirement element}, a formally identified statement of what the research must measure. It models each survey item as a \emph{Block}, the basic building-block unit the model uses to represent one distinct question. It then binds every item to the requirement it addresses through a \emph{Satisfy} relationship, the modeling link that records that one element fulfills another. These links form a \emph{digital thread}, a connected and authoritative chain of references that ties each item to its research question. Every item therefore traces, explicitly and bidirectionally, to the research question it is intended to measure. Completeness of that mapping can be evaluated programmatically rather than asserted in prose. The approach supplies documented, quantitative traceability and content-validity evidence in a form that is at once visually inspectable (Figures~\ref{fig:MBSERQ} and~\ref{fig:ResearchQuestionTraceability}) and machine-checkable. It does so by applying to the research instrument the very digital-thread capability whose value the dissertation sets out to assess. The method thus demonstrates in practice the traceability and rigor it investigates rather than merely describing them.

The survey model follows a structured lifecycle consisting of six phases:

\subsubsection{Strategic Phase}

Captures conceptual elements including hypothesis, capabilities, constraints, goals, vision, timeline, milestones, stakeholders, and drivers. This phase establishes the foundation for all subsequent research activities by clearly articulating what the research seeks to accomplish and why it matters.

\subsubsection{Requirements Phase}

Derives and analyzes strategic elements to create requirements for the dissertation and associated survey. Requirements follow ISO 15288:2023 standards: hierarchical in nature and structure, atomic in composition, bidirectional in traceability, individually measurable, and explicitly tracking relationships and interdependencies.

\subsubsection{Architecture Phase}

Establishes the high-level outline, structure, and guidance for the design and implementation phases of the dissertation. The survey architecture defines the overall structure of sections and question flow, ensuring logical progression from awareness through applicability to perceived value.

\subsubsection{Design Phase}

The Design Phase uses the architecture guidance and requirements to design the dissertation proposal and survey instrument, satisfying the architectural outline with complete traceability back to requirements and strategic elements.

\subsubsection{Results Phase}

Captures the survey data into the model for traceability along with results of analysis. Analysis methods and results are captured as elements within the model, enabling verification that findings address the original research questions.

\subsubsection{Report Phase}

The model serves as a deliverable component of the dissertation that provides a rigorous source of artifacts, methods used, traceability chains, and evidentiary support for the findings.

\subsubsection{Alignment to Dissertation Design}

Figure~\ref{fig:DissertationLifecycleProcess} provides a visual representation of the dissertation process using an MBSE approach. The research has progressed through the full systems engineering lifecycle: the strategic, requirements, architecture, and design phases established the survey instrument and analytical plan, while the results and report phases (reported in Chapter~4) captured the collected survey data and the results of analysis into the model for traceability.

\begin{figure}[ht]
  \centering
  \includegraphics[width=0.8\linewidth, alt={Abstract resource taxonomy of the dissertation lifecycle as a vertical chain of the six systems-engineering phases the research followed: Strategic, Requirements, Architecture, Design, Results, and Report, each phase producing an output that feeds the next: Strategic Elements, Requirements, Architecture Guidance, Survey Instrument, and Analysis Results.}]{figures/DissertationAbstractResourceTaxonomy.pdf}
  \caption{Dissertation lifecycle process as an abstract resource taxonomy.}
  \label{fig:DissertationLifecycleProcess}
\end{figure}

\clearpage

\subsection{Research Questions and Survey Alignment}

The survey instrument addresses the three research questions directly:

\subsubsection{Research Question 1}

To what extent are Information Technology and Information Assurance professionals aware of Digital Engineering capabilities, including Model-Based Systems Engineering, digital threads, digital twin technologies, and Product Lifecycle Management principles?

\subsubsection{Research Question 2}

To what extent do Information Technology and Information Assurance professionals perceive Digital Engineering capabilities as potentially valuable or important for their work in Information Assurance, security compliance, and IT service delivery?

\subsubsection{Research Question 3}

To what extent do Information Technology and Information Assurance professionals believe that Digital Engineering practices could help them in performing their jobs, meeting compliance requirements, or enhancing organizational capabilities in Information Assurance and IT service delivery?

\subsubsection{Research Propositions and Planned Analyses}

Chapter 1 advances four research propositions that follow from the theoretical framework; Chapter 4 evaluates each proposition against the survey data, and Chapter 5 returns to them in stating the conclusions. The analytical plan binds each proposition to a designated analysis:

\begin{itemize}
  \item \textbf{P1 (Low awareness).} Awareness of and recent exposure to Digital Engineering are low among IT and IA professionals, who sit outside the systems engineering discipline in which these methods developed. Descriptive statistics on the familiarity (Question 1.1) and exposure (Question 1.2) items, reported with Wilson confidence intervals, assess this proposition.
  \item \textbf{P2 (Value despite low awareness).} Perceived value is nonetheless positive once capabilities are described in functional terms, because Digital Engineering addresses the documentation, traceability, and visibility problems these professionals routinely face. Descriptive agreement levels on the value items, together with the familiarity-versus-value diagnostic described in the Response Bias Mitigation section, assess this proposition.
  \item \textbf{P3 (Variation by field and experience).} Perceived value and awareness vary systematically with professional field and experience level rather than being uniform across the community. Mann--Whitney U tests (professional field), Kruskal--Wallis tests with Dunn's post-hoc comparisons under Holm adjustment (experience level), and chi-square tests with Cram\'er's V (categorical associations) assess this proposition.
  \item \textbf{P4 (Familiarity and value).} Greater familiarity with Digital Engineering is associated with higher perceived value. Rank-order correlation and ordinary least squares regression of the value composites on familiarity, professional field, and experience level assess this proposition.
\end{itemize}

The propositions above define the analyses specified in this proposal. Chapter~4 additionally reports a set of supplementary robustness and measurement-validity analyses that this plan did not specify, including a common-method-variance screen, a construct-validity assessment of the composite scores, and a measurement-invariance check across professional-field groups. The candidate added those analyses after data collection concluded, as a self-initiated methodological strengthening informed by comparison with related dissertations; Chapter~4 labels them as post-hoc where they appear and interprets them as robustness evidence rather than as tests of the propositions.

The design mapped each survey section and question to the three research questions to support comprehensive coverage of the research objectives while minimizing respondent burden through focused questioning. This mapping provides the traceability essential to demonstrating that the instrument measures what it purports to measure.

Figure~\ref{fig:MBSERQ} illustrates an MBSE representation of the Research Questions as requirement type elements. Selecting requirement type elements within the Unified Architecture Framework (UAF) allows survey questions to use the ``Satisfy'' relationship. Authoritative traceability within the model enables programmatic evaluation that verifies the completeness of the question-to-requirement mapping; this check validates traceability coverage rather than construct validity.

\clearpage

\begin{landscape}
  \begin{figure}[H]
    \centering
    \includegraphics[width=1\linewidth, alt={SysML requirement diagram showing the three research questions as requirement elements, each with an identifier and text, linked by Satisfy dependencies to the survey-question Blocks that measure it; each Block carries a Likert, Binary, or Ternary response operation.}]{figures/ResearchQuestionTaxonomy.pdf}
    \caption{MBSE instantiation of research questions.}
    \label{fig:MBSERQ}
  \end{figure}
\end{landscape}

\clearpage

\section{Survey Requirements Specification}

The systems engineering approach established specific instrument requirements before design. These requirements help align the survey with the research objectives while maintaining ethical standards and data quality.

\subsection{Anonymity and Privacy Requirements}

The following requirements govern anonymity and privacy protection:

\begin{itemize}
  \item The survey shall be anonymous, collecting no personally identifiable information during the survey process.
  \item The survey shall not collect age, sex, gender, race, creed, gender identity, sexual orientation, income, nationality, or other non-relevant personal information.
  \item The survey shall be protected with appropriate security controls to ensure the confidentiality, integrity, and availability of collected data.
  \item The survey raw data shall not be sold or used for any purpose other than the dissertation research.
\end{itemize}

\subsection{Instrument Design Requirements}

The following requirements govern survey instrument design:

\begin{itemize}
  \item The survey shall use a Likert Scale for material research questions to enable quantitative analysis.
  \item The survey shall include applicable notices required by the University Institutional Review Board.
  \item The survey should collect demographic information to identify the participant's professional field of study or practice.
  \item The survey should collect demographic information about temporal experience level within the participant's field.
\end{itemize}

\subsection{Content Requirements}

The survey content assesses the awareness, applicability, and perceived value that participants hold toward Digital Engineering and its component disciplines. The survey begins at the enterprise level and examines the four pillars: Model-Based Systems Engineering, the Digital Thread, Digital Twin, and Product Lifecycle Management. Consistent question flow across this structure enables analysis of how each discipline is perceived in relation to Digital Engineering as a unified approach.

Value assessment questions in Sections 3 through 5 use positively framed capability statements (e.g., ``Digital Engineering could deliver meaningful value'' or ``could reduce development cycle time''). The research objectives and target population characteristics informed this framing through deliberate design decisions. Because the research questions ask to what extent professionals perceive value in Digital Engineering capabilities, positively framed statements directly measure this construct by presenting capabilities and assessing agreement. An instrument measuring perceived value necessarily presents value propositions for evaluation. Negatively framed items (e.g., ``Digital Engineering would not improve security posture'') would introduce cognitive complexity through double negatives when combined with disagreement responses. The survey methodology literature documents this confound, which disproportionately affects respondents unfamiliar with the subject matter.

The design addresses the acquiescence bias risk inherent in consistently framed questions through two analytical checks. First, response-pattern analysis examined whether any respondents showed uniform agreement across all items (straight-lining), a pattern suggestive of satisficing or acquiescence rather than considered evaluation. Second, the Section 1 familiarity self-assessment provides an independent baseline against which value perception scores can be interpreted. Respondents who report low familiarity yet indicate high agreement with value statements can be identified and examined. Such a pattern signals potential acquiescence rather than genuine perception. Two response-format features complement these checks. A neutral midpoint (``Neither agree nor disagree'') reduces forced-choice artifacts by providing a non-affirmative response option for respondents who are genuinely uncertain or indifferent. The binary and ternary investment willingness questions (Questions 4.5 and 5.7) offer explicit ``No'' and ``Unsure'' options. These options capture skepticism and uncertainty without requiring disagreement with a positively framed statement. Because all value items are positively keyed, these safeguards detect and bound acquiescence rather than eliminate it, and residual acquiescence risk remains an acknowledged limitation of the instrument.

\section{Target Population and Sampling}

\subsection{Target Population}

Professionals actively working in Information Technology and Information Assurance roles within organizations that manage information systems constitute the target population for this study. The population encompasses individuals involved in IT service delivery, infrastructure management, service operations, IA, security operations, compliance management, and security architecture. It spans organizational types including private sector enterprises, government agencies, educational institutions, and nonprofit organizations.

The dual focus on IT and IA professionals recognizes that Digital Engineering capabilities may offer distinct value propositions for these related but different professional communities. IT professionals primarily concerned with service delivery and operational efficiency may perceive value in Digital Engineering practices differently than IA professionals focused on threat mitigation, compliance verification, and security control implementation. Understanding both perspectives is necessary for thorough assessment.

\subsection{Sampling Strategy}

This study uses non-probability convenience sampling, recruiting participants through professional networking platforms and direct professional outreach. While probability sampling would provide stronger generalizability, establishing a sampling frame for all IT and Information Assurance professionals is difficult, which renders probability sampling impractical for this research. Convenience sampling enables efficient access to qualified respondents while maintaining focus on professionals with relevant experience and expertise.

The design anticipated recruiting across several channels, among them professional organizations such as ISACA, (ISC)\textsuperscript{2}, and ITIL professional communities; LinkedIn professional groups focused on IT service management and IA; industry conferences and professional development events; and direct outreach to IT and security departments across multiple sectors. Execution did not realize that breadth. As the Recruitment Channels and Promotion section below records, LinkedIn served as the primary medium, and because the candidate's network there is concentrated in defense, aerospace, and federal organizations, the realized frame is narrower than the plan and the multi-channel breadth cannot be claimed as a mitigation. Chapter~5 carries the resulting composition concern into the limitations, where it is the most consequential the study reports.

\subsection{Sample Size Determination}

The sample size calculation targeted a maximum margin of error of 5\% at the 95\% confidence level. For survey research targeting a large population where the population size exceeds 100,000 individuals, the sample size formula for estimating a proportion is:
\clearpage
\begin{equation}
  \label{eqn:modelSizeEquation}
  n = \frac{Z^2 \times p \times (1-p)}{E^2}
\end{equation}
\addcontentsline{lofm}{formula}{\protect\numberline{\theequation}Sample size for estimating a proportion}

Where:
\begin{itemize}
  \item \(n\) = required sample size
  \item \(Z\) = Z-score for desired confidence level (1.96 for 95\% confidence)
  \item \(p\) = estimated population proportion (0.5 for maximum variability)
  \item \(E\) = desired margin of error (0.05 for 5\%)
\end{itemize}

Substituting these values:

\begin{equation}
  \label{eqn:ModelSizeCalc}
  n = \frac{1.96^2 \times 0.5 \times 0.5}{0.05^2} = \frac{3.8416 \times 0.25}{0.0025} = \frac{0.9604}{0.0025} = 384.16
\end{equation}
\addcontentsline{lofm}{formula}{\protect\numberline{\theequation}Sample size with substituted values}

Therefore, a minimum of 385 completed survey responses is required to achieve the target margin of error. This calculation assumes maximum variability in responses (p = 0.5), which provides the most conservative sample size estimate. If actual response distributions show less variability, the effective margin of error will be smaller than 5\%. These precision calculations assume probability sampling; under the convenience sampling used here they serve as planning benchmarks, and reported confidence intervals are conditional on the sample behaving like a random draw from the accessible population.

To account for potential incomplete responses and to support adequate representation across professional subgroups (IT professionals versus security professionals, varying experience levels), the study set the target sample size for data collection at 450 completed surveys. This target exceeds the minimum requirement by approximately 17\%, aligning with survey research best practices recommending 15\% to 20\% oversampling to accommodate anticipated attrition~\parencite{Dillman_Smyth_Christian_2014}. Research on online survey completion rates indicates that professional surveys typically experience 10\% to 15\% incomplete response rates due to survey abandonment, partial completion, or data quality concerns requiring exclusion~\parencite{Fan_Yan_2010}. The 450-response target provides sufficient buffer to reach the minimum 385 usable responses while enabling meaningful subgroup analyses with adequate statistical power. Specifically, if the sample divides evenly between IT and security professionals, each subgroup would include approximately 190--225 respondents, providing margins of error of roughly 6.5\% to 7.1\% for subgroup-specific analyses. The realized sample substantially exceeded this target. The survey drew 1{,}600 submissions. The data-quality screening detailed in the Data Analysis section then removed 189 cases (188 abandoned or aberrant submissions, plus one automated entry identified by the response-authenticity check reported in Chapter~4), approximately 12\%, a figure squarely within the attrition range anticipated above. After that screening, 1{,}411 responses remained for analysis. The retained sample, reported in Chapter~4, yields subgroup sizes well above the planning assumptions and correspondingly tighter precision than the design required.

\section{Survey Instrument Design}

\subsection{Instrument Overview}

The 30-question survey instrument systematically addresses the research questions across six thematic sections while maintaining reasonable completion time. Question complexity and respondent fatigue considerations set the target completion time at roughly ten minutes, balancing comprehensive data collection against respondent engagement and completion rates.

The six sections are structured as follows:

\begin{enumerate}
  \item Section 1: Awareness and Familiarity with Digital Engineering (2 questions)
  \item Section 2: Understanding of Digital Engineering Capabilities (6 questions)
  \item Section 3: Applicability of Digital Engineering (6 questions)
  \item Section 4: Value Assessment for Information Technology (5 questions)
  \item Section 5: Value Assessment for Information Assurance and Cybersecurity (7 questions)
  \item Section 6: Interest and Demographic Information (4 questions)
\end{enumerate}

\subsection{Question Format and Scale Selection}

The survey uses two primary question formats: five-point Likert-type scales and binary (yes/no) response options. Psychometric research establishing the validity and reliability of these formats for measuring attitudes, perceptions, and self-reported behaviors guided this choice.

\subsubsection{Likert Scale Justification}

Five-point Likert scales measure awareness levels, agreement with capability statements, and perceptions of value. Likert scales have been extensively validated for measuring attitudes and perceptions in social science research since their introduction by Rensis Likert in 1932. Research showing that five to seven response options provide optimal discrimination between respondent positions while remaining cognitively manageable drove the choice of five options.

The analytical treatment of Likert scale data warrants explicit methodological comment. A longstanding debate in psychometric and social science research concerns whether Likert scale responses constitute ordinal data (where distances between scale points are not necessarily equal) or may be treated as interval data (where parametric statistical methods such as means and standard deviations are appropriate). Jamieson~\parencite{Jamieson_2004} argues that Likert data are strictly ordinal and that parametric analyses may produce misleading results. Norman~\parencite{Norman_2010}, drawing on extensive simulation research, shows that parametric methods are sufficiently robust to yield valid results with Likert data even when distributional assumptions are violated. Subsequent research by Sullivan and Artino~\parencite{Sullivan_Artino_2013} and de Winter and Dodou~\parencite{deWinter_Dodou_2010} supports Norman's position, finding that parametric tests applied to five-point Likert data produce conclusions consistent with non-parametric alternatives across a wide range of distributional conditions.

This research follows the analytical convention adopted by most technology acceptance and survey research. It reports parametric descriptive statistics (means, standard deviations) for Likert items and composite scores to enable comparison with the broader literature. When distributional assumptions are not met, it uses non-parametric alternatives for inferential testing. Specifically, the distributional shape of each composite (skewness and proximity to the response ceiling), corroborated by Shapiro--Wilk tests and visual inspection, determines whether parametric tests (t-tests, ANOVA) or non-parametric alternatives (Mann--Whitney U, Kruskal--Wallis) are applied for group comparisons. The Shapiro--Wilk test is treated as corroborating rather than dispositive, because at the realized sample size it is highly powered against trivial departures from normality and would reject for substantively negligible deviations; the decision therefore rests primarily on the observed left skew and ceiling concentration of the composites. The analysis reports medians and interquartile ranges alongside means and standard deviations for individual Likert items, providing readers with both parametric and ordinal summaries. This dual-reporting approach supports transparency and enables readers who hold either position in the ordinal-interval debate to evaluate the findings against their preferred analytical framework.

Consistent use of Likert scales across sections provides an intuitive response method for participants and enables direct comparison of responses across different constructs. Two distinct anchor sets serve the constructs being measured:

\paragraph{Familiarity Scale (Section 1, Question 1.1)}
\begin{enumerate}
  \item Not at all familiar
  \item Slightly familiar
  \item Moderately familiar
  \item Very familiar
  \item Extremely familiar
\end{enumerate}

\paragraph{Agreement Scale (Sections 2-5)}
\begin{enumerate}
  \item Strongly disagree
  \item Disagree
  \item Neither agree nor disagree
  \item Agree
  \item Strongly agree
\end{enumerate}

The agreement scale follows established conventions in technology acceptance research, including the Technology Acceptance Model (TAM) and Unified Theory of Acceptance and Use of Technology (UTAUT) frameworks, and has demonstrated validity for measuring perceptions of technology value and usefulness across diverse populations and contexts. Including a neutral midpoint (``Neither agree nor disagree'') enables respondents to indicate genuine uncertainty rather than forcing artificial choices, which research suggests improves response validity. Figure~\ref{fig:SurveyEnumerations} shows the instantiation of familiarity and agreement response scales within the model of the dissertation survey.

\subsubsection{Binary and Ternary Question Justification}

The survey uses binary yes/no and ternary yes/no/unsure formats for factual questions about prior exposure (Section 1, Question 1.2), interest in further learning (Section 6, Questions 6.1-6.2), and investment willingness (Sections 4-5, Questions 4.5 and 5.7). Binary/ternary formats are appropriate for these questions because they seek to classify respondents into discrete categories rather than measure degrees of agreement or perception intensity. The simplicity of binary/ternary responses also reduces cognitive load for straightforward factual questions where graduated responses would add complexity without meaningful information gain. Figure~\ref{fig:SurveyEnumerations} provides a visualization of the enumeration elements of the binary and ternary responses.

\subsubsection{Categorical Questions}

Section 6 includes categorical questions capturing demographic information about professional field of practice and years of experience. Following the demographic information requirements, these questions capture only sector of work and temporal experience level without collecting personally identifiable information. The experience level categories are structured to capture early-career, mid-career, and late-career stages:

\begin{itemize}
  \item 1-5 Years of experience (Early Career)
  \item 6-10 Years of experience (Mid Career - Early)
  \item 11-15 Years of experience (Mid Career - Late)
  \item 16+ Years of experience (Late Career)
\end{itemize}

These categories enable analysis of response patterns across professional subgroups and experience levels, supporting investigation of whether awareness and perceptions vary systematically by professional background. Such variation would carry important implications for professional development and adoption strategies.

\begin{figure}[H]
  \centering
  \includegraphics[width=1\linewidth, alt={SysML package diagram of the six survey response-scale enumerations: PracticeExperience (1-5, 6-10, 11-15, and 16+ years), PracticeField (IT, Information Assurance/Cyber Security, Engineering, Other), BinaryAffirmation (Yes, No), TernaryAffirmation (Yes, No, Unsure), LikertFamiliarity (five levels from Not at all familiar to Extremely familiar), and LikertAgreement (five levels from Strongly disagree to Strongly agree).}]{figures/OverviewDiagramOfPackageEnumerations.pdf}
  \caption{Survey question enumeration.}
  \label{fig:SurveyEnumerations}
\end{figure}

\section{Survey Section Structure and Research Question Mapping}

This section details how each survey section addresses the research questions, providing explicit traceability between survey content and research objectives. Every question serves a defined purpose, and the mapping supports comprehensive coverage of the research questions.

Table~\ref{tab:survey-items} consolidates all thirty items. For each item it records the construct measured, the verbatim item stem, and the research question or category to which the item maps. The subsections that follow describe each section's role and its mapping to the research questions.

{\small
\setlength{\tabcolsep}{4pt}
\begin{longtable}{@{}l p{2.9cm} p{7.5cm} l@{}}
\caption{Survey items mapped to measured constructs and research questions.}\label{tab:survey-items}\\
% ALT: Complete list of the 30 survey items (Q1.1 through Q6.4). Each row gives the item number, the construct it measures, the verbatim item stem, and the mapped research question or category. Section 1 and Section 2 items map to Research Question 1; Section 3, 4, and 5 items map to Research Questions 2 and 3; Section 6 items map to Supporting and Demographic categories.
\toprule
Item & Construct measured & Verbatim item stem & Mapped RQ \\
\midrule
\endfirsthead
\multicolumn{4}{@{}l}{\small\itshape Table~\ref{tab:survey-items} continued from previous page}\\
\toprule
Item & Construct measured & Verbatim item stem & Mapped RQ \\
\midrule
\endhead
\midrule
\multicolumn{4}{r@{}}{\small\itshape Continued on next page}\\
\endfoot
\bottomrule
\endlastfoot
Q1.1 & Self-reported awareness levels & ``Please rate your level of familiarity with Digital Engineering concepts and practices.'' & RQ1 \\
Q1.2 & Recent professional exposure & ``Have you encountered Digital Engineering methodologies, frameworks, or tools in your professional work within the past two years?'' & RQ1 \\
Q2.1 & Model-Based Systems Engineering & ``Digital Engineering includes model-based systems engineering approaches that can improve development processes.'' & RQ1 \\
Q2.2 & Digital Twin for IT & ``Digital Engineering can enable digital twin development and virtual prototyping for Information Technology systems.'' & RQ1 \\
Q2.3 & Digital Thread and data-driven practices & ``Digital Engineering supports continuous integration and data-driven decision-making in technology development.'' & RQ1 \\
Q2.4 & Digital Twin for security & ``Digital Engineering enables digital twin technology that can simulate security scenarios and test defensive measures without impacting production systems.'' & RQ1 \\
Q2.5 & Security validation through the Digital Thread & ``Digital Engineering supports continuous security validation and data-driven threat analysis throughout the development lifecycle.'' & RQ1 \\
Q2.6 & Compliance automation & ``Digital Engineering can improve security control implementation through automated compliance checking and verification.'' & RQ1 \\
Q3.1 & IT sector relevance & ``Digital Engineering methodologies have relevant applications within the Information Technology sector.'' & RQ2 \\
Q3.2 & Information Assurance relevance & ``Digital Engineering methodologies have relevant applications for addressing Information Assurance challenges.'' & RQ2 \\
Q3.3 & Digital Twin organizational value & ``The ability to utilize digital twins to test changes against accurate replicas of production environments would provide value to my organization.'' & RQ2, RQ3 \\
Q3.4 & MBSE organizational value & ``The use of digital models to map and document an organization's environment and configurations would provide value to my organization.'' & RQ2, RQ3 \\
Q3.5 & PLM organizational value & ``The use of digital lifecycle management to meet compliance and service delivery requirements would provide value to my organization.'' & RQ2, RQ3 \\
Q3.6 & Compliance applicability & ``My organization faces regulatory or compliance requirements that could benefit from Digital Engineering approaches.'' & RQ2, RQ3 \\
Q4.1 & Overall IT value & ``Digital Engineering could deliver meaningful value to my organization's Information Technology processes.'' & RQ2, RQ3 \\
Q4.2 & Cycle time benefit & ``Digital Engineering could reduce development cycle time in my organization.'' & RQ3 \\
Q4.3 & Quality benefit & ``Digital Engineering could improve product quality and reduce defects in my organization.'' & RQ3 \\
Q4.4 & Collaboration benefit & ``Digital Engineering could improve collaboration effectiveness across development teams in my organization.'' & RQ3 \\
Q4.5 & Investment willingness & ``My organization would be willing to invest in Digital Engineering capabilities if clear return on investment could be demonstrated.'' (Yes/No/Unsure) & RQ2 \\
Q5.1 & Overall security value & ``Digital Engineering could deliver meaningful value to my organization's Information Assurance and cybersecurity operations.'' & RQ2, RQ3 \\
Q5.2 & Vulnerability management benefit & ``Digital Engineering could reduce the time required to identify and remediate security vulnerabilities in my organization.'' & RQ3 \\
Q5.3 & Security posture benefit & ``Digital Engineering could improve security posture and reduce successful cyber incidents in my organization.'' & RQ3 \\
Q5.4 & Threat modeling benefit & ``Digital Engineering could enhance threat modeling and risk assessment capabilities in my organization.'' & RQ3 \\
Q5.5 & Cross-team collaboration benefit & ``Digital Engineering could improve collaboration between security teams, development teams, and operations teams in my organization.'' & RQ3 \\
Q5.6 & Compliance benefit & ``Digital Engineering could help my organization achieve better compliance with security frameworks and regulatory requirements.'' & RQ3 \\
Q5.7 & Security investment willingness & ``My organization would be willing to invest in Digital Engineering capabilities for cybersecurity purposes if clear return on investment could be demonstrated.'' (Yes/No/Unsure) & RQ2 \\
Q6.1 & Learning interest & ``Would you be interested in learning more about Digital Engineering applications for Information Assurance, Cybersecurity, or Information Technology in your industry?'' & Supporting \\
Q6.2 & Recommendation likelihood & ``Would you recommend that your organization explore Digital Engineering adoption for improving security operations?'' & Supporting \\
Q6.3 & Professional field & ``Please indicate your field of practice.'' (Information Technology / Information Assurance or Cyber Security / Engineering / Other) & Demographics \\
Q6.4 & Experience level & ``Please indicate your level of experience in your field of practice.'' (1-5 Years / 6-10 Years / 11-15 Years / 16+ Years) & Demographics \\
\end{longtable}
}

\subsection{Section 1: Awareness and Familiarity with Digital Engineering}

Section 1 contains two questions directly addressing Research Question 1 regarding professional awareness of Digital Engineering capabilities. Figure~\ref{fig:Section1Questions} provides an MBSE-based visualization as a block element, and Figure~\ref{fig:SurveyQuestionsDetailedProperties} details the properties recorded for each question element. Table~\ref{tab:survey-items} lists both item stems. Question 1.1 uses a five-point familiarity scale and provides a direct, self-assessed measure of overall Digital Engineering awareness, establishing the baseline from which all subsequent analysis proceeds. Question 1.2 uses a binary format and distinguishes theoretical awareness from practical professional experience, recognizing that knowing about Digital Engineering differs categorically from having encountered it in practice.

\begin{figure}[H]
  \centering
  \includegraphics[width=1\linewidth, alt={SysML diagram detailing the recorded properties of the Section 1 survey-question elements, including question text and response type.}]{figures/OverviewDiagramOfPackageSection1DetailedProperties.pdf}
  \caption{Survey questions Section 1 detailed properties taxonomy.}
  \label{fig:SurveyQuestionsDetailedProperties}
\end{figure}

\begin{figure}[H]
  \centering
  \includegraphics[width=1\linewidth, alt={SysML package diagram for survey Section 1, containing two Block elements: Question 1.1, a Likert familiarity item, and Question 1.2, a binary item.}]{figures/OverviewDiagramOfPackageSection1.pdf}
  \caption{Section One survey question elements package diagram.}
  \label{fig:Section1Questions}
\end{figure}

\subsubsection{Mapping to Research Question 1}

These questions directly measure the extent of professional awareness, providing quantifiable data on familiarity levels and recent professional exposure to Digital Engineering concepts.

\subsection{Section 2: Recognition of Digital Engineering Capabilities}

Section 2 contains six Likert-scale questions assessing respondent recognition of specific Digital Engineering capabilities, each described in functional terms. Figure~\ref{fig:Section2Questions} visualizes each research question as represented in the model. The questions of Section 2 probe awareness of the four Digital Engineering pillars in the context of IT and IA applications, as listed in Table~\ref{tab:survey-items}.

\begin{figure}[H]
  \centering
  \includegraphics[width=1\linewidth, alt={SysML package diagram for survey Section 2, containing six Block elements, Questions 2.1 through 2.6, each a Likert agreement item.}]{figures/OverviewDiagramOfPackageSection2.pdf}
  \caption{Section Two survey question elements package diagram.}
  \label{fig:Section2Questions}
\end{figure}

\subsubsection{Mapping to Research Question 1}

This section extends Research Question 1 by measuring both general awareness and recognition of specific Digital Engineering pillars (MBSE, Digital Twin, the Digital Thread, and PLM) and their applications to IT and IA contexts. A methodological distinction warrants clarification: Section 2 questions use agreement scales rather than direct awareness or knowledge measures. Agreement with a capability statement (e.g., ``Digital Engineering includes model-based systems engineering approaches that can improve development processes'') captures recognition of a described capability, meaning whether respondents affirm the stated relationship once it is put to them in functional language. Establishing comprehension would require knowledge-verification items, such as definitional or discrimination items carrying a correct answer, and the instrument includes none. Section 1 captures familiarity with the named discipline directly through self-assessment; Section 2 probes further by testing whether respondents affirm the capability descriptions that give the discipline its content. Taken together, Sections 1 and 2 provide a layered assessment of Research Question 1 that separates familiarity with the label from recognition of the capabilities described functionally, and the design claims no measure of knowledge beyond those two.

\subsection{Section 3: Applicability of Digital Engineering}

Section 3 contains six Likert-scale questions examining perceptions of Digital Engineering applicability to IT and IA domains. This section bridges awareness assessment and value perception, probing whether respondents see relevance to their professional contexts. Figure~\ref{fig:Section3Questions} visualizes the instantiation of Section 3 survey questions within the model.

\begin{figure}[H]
  \centering
  \includegraphics[width=1\linewidth, alt={SysML package diagram for survey Section 3, containing six Block elements, Questions 3.1 through 3.6, each a Likert agreement item.}]{figures/OverviewDiagramOfPackageSection3.pdf}
  \caption{Section Three survey question elements package diagram.}
  \label{fig:Section3Questions}
\end{figure}

\subsubsection{Mapping to Research Questions 2 and 3}

This section addresses whether respondents perceive Digital Engineering as applicable to their professional domains and whether they identify potential organizational value in specific capabilities.

\subsection{Section 4: Value Assessment for Information Technology}

Section 4 contains five questions assessing the perceived value of Digital Engineering for IT operations specifically; Figure~\ref{fig:Section4Questions} visualizes these questions.

\begin{figure}[H]
  \centering
  \includegraphics[width=1\linewidth, alt={SysML package diagram for survey Section 4, containing five Block elements: Questions 4.1 through 4.4 as Likert agreement items and Question 4.5 as a ternary item.}]{figures/OverviewDiagramOfPackageSection4.pdf}
  \caption{Section Four survey question elements package diagram.}
  \label{fig:Section4Questions}
\end{figure}

\subsubsection{Mapping to Research Questions 2 and 3}

This section directly addresses whether IT professionals perceive value in Digital Engineering and believe these practices could enhance their organizational capabilities and job performance.

\subsection{Section 5: Value Assessment for Information Assurance and Cybersecurity}

Section 5 contains seven questions assessing perceived value of Digital Engineering for Information Assurance and cybersecurity operations as displayed in Figure~\ref{fig:Section5Questions}.

\begin{figure}[H]
  \centering
  \includegraphics[width=1\linewidth, alt={SysML package diagram for survey Section 5, containing seven Block elements: Questions 5.1 through 5.6 as Likert agreement items and Question 5.7 as a ternary item.}]{figures/OverviewDiagramOfPackageSection5.pdf}
  \caption{Section Five survey question elements package diagram.}
  \label{fig:Section5Questions}
\end{figure}

\subsubsection{Mapping to Research Questions 2 and 3}

This section directly addresses whether IA professionals perceive value in Digital Engineering for their specific domain and believe these practices could help them meet compliance requirements and enhance security capabilities.

\subsection{Section 6: Interest and Demographic Information}

Section 6 contains four questions capturing respondent interest in Digital Engineering learning opportunities and demographic characteristics. Figure~\ref{fig:Section6Questions} visualizes Section 6 survey questions.

\subsubsection{Question 6.3}

Question 6.3 captures the respondent's professional field; Table~\ref{tab:survey-items} gives its verbatim stem and response options. Combining Information Assurance and Cyber Security into a single response category reflects a practical design choice: working professionals identify their field using varied terminology including cybersecurity, information security, and Information Assurance, and a single combined label maximizes correct self-selection without presuming which term a respondent uses. The category was worded to read as familiar both to the governance-oriented professionals Chapter~1 designates Information Assurance and to the operationally oriented professionals it designates cybersecurity. One consequence, carried into the limitations of Chapter~5, is that the resulting category is internally heterogeneous, spanning both senses, and that the anonymous instrument cannot separate a governance-oriented from an operationally oriented respondent within it after the fact. Respondents selecting Engineering or Other appear in descriptive analyses; the primary field comparison contrasts the Information Technology and Information Assurance/Cyber Security groups.

\begin{figure}[H]
  \centering
  \includegraphics[width=1\linewidth, alt={SysML package diagram for survey Section 6, containing four Block elements: Questions 6.1 and 6.2 as binary items and Questions 6.3 and 6.4 capturing professional field and experience.}]{figures/OverviewDiagramOfPackageSection6.pdf}
  \caption{Section Six survey question elements package diagram.}
  \label{fig:Section6Questions}
\end{figure}

\subsubsection{Mapping to Research Questions}

While not directly measuring awareness or perceived value, demographic data enables investigation of whether awareness and perceptions vary by professional field or experience level, providing insights for targeted future research and professional development initiatives.

\subsection{Traceability of Survey Questions}

Figure~\ref{fig:ResearchQuestionTraceability} provides a comprehensive diagram that shows authoritative traceability of survey question elements to the research questions. Relationship traceability diagrams enable programmatic verification that every survey question traces to a research question; this check validates traceability completeness rather than construct validity. The mapping is complete: every survey item resolves to at least one research question or supporting demographic category, and no item stands outside the traceability structure.

\begin{figure}[H]
  \centering
  \includegraphics[width=1\linewidth, alt={Satisfy traceability matrix mapping the five research-question groups in rows (Research Questions 1, 2, and 3, Demographics, and Supporting) to all survey questions across Sections 1 through 6 in columns, with arrows marking which questions satisfy each group; the row totals are 8, 10, 14, 2, and 2 respectively.}]{figures/ResearchQuestionTraceabilityMatrix.pdf}
  \caption{Research question traceability.}
  \label{fig:ResearchQuestionTraceability}
\end{figure}

\section{Data Collection Procedures}

\subsection{Survey Platform and Administration}

The study administered the survey electronically through Microsoft Forms, which supports anonymous response collection, secure data storage, and compliance with research ethics requirements. % Confirmed by author 2026-07-02: platform was Microsoft Forms.
Electronic administration enabled efficient access to geographically distributed respondents while ensuring consistent question presentation and response capture across all participants. Collection ran from April 23 through June 2026, opening once the Institutional Review Board determination was in place and closing when the realized sample had passed the target described above.

\subsection{Data Protection Plan}

In accordance with the survey requirements, the study protected the survey data through the specific controls listed below, safeguarding the confidentiality, integrity, and availability of collected data. The data protection plan included:

\begin{itemize}
  \item Password protection and multifactor authentication for access to survey administration and data
  \item Encryption at rest for stored survey responses
  \item Encryption in transit for data transmission between respondents and the survey platform
  \item No collection of IP addresses or other metadata that could potentially identify respondents
  \item Secure storage in compliance with university data governance requirements
\end{itemize}

\subsection{Informed Consent}

The survey begins with a comprehensive participation notice that explains the research purpose, what participation involves, confidentiality protections, the voluntary nature of participation, and contact information for questions or concerns. Specifically, the notice states:

\begin{itemize}
  \item The survey is conducted as part of a doctoral research project at Dakota State University
  \item Participation is completely anonymous with no personally identifiable information collected
  \item Risks of participation are minimal and similar to those encountered in normal daily internet activity
  \item Participation is entirely voluntary with the right to withdraw at any time before submission
  \item Respondents must be 18 years or older to participate
  \item No compensation or incentive is offered for participation
  \item Completion and submission of the survey constitutes informed consent
\end{itemize}

\subsection{Recruitment and Distribution}

The candidate distributed the survey through professional networking outreach, detailed by channel in the following subsection. Recruitment messages clearly described the target population (IT and IA professionals), estimated completion time (approximately ten minutes), and research purpose. The announcement was framed descriptively rather than persuasively, naming the study topic and eligible population without asserting that Digital Engineering delivers value, and it is reproduced verbatim in Appendix~A.

\subsection{Recruitment Channels and Promotion}

LinkedIn served as the primary recruitment medium. Direct participation requests went to professionals in the candidate's LinkedIn network, without filtering by industry or employer. Network composition rather than a sampling decision therefore determined which sectors the invitation reached, and those contacts span aerospace, defense, technology, and government organizations whose work intersects with systems engineering, IT, and IA. These organizations included Collins Aerospace, RTX, Raytheon Technologies, Pratt \& Whitney, Anduril, General Atomics, General Dynamics Mission Systems, Boeing, Lockheed Martin, Northrop Grumman, Shield AI, Booz Allen Hamilton, Broadcom, Palo Alto Networks, Dell, Supermicro, Microsoft, the Defense Information Systems Agency, the National Security Agency, the Department of War, Cisco, Defense Unicorns, General Electric, and INCOSE. Reach therefore spanned both the defense and aerospace sector in which Digital Engineering originated and the broader enterprise IT and IA community the research seeks to reach, in proportions the network set rather than the design. Because that network is itself concentrated in defense, aerospace, and federal organizations, the realized sample is weighted toward those sectors, a composition whose implications for generalizability are addressed in the limitations of Chapter~5.

To widen the reach of the recruitment post beyond the candidate's immediate professional network, a paid LinkedIn promotion campaign ran from early May through mid-June with a capped spend of 500 United States dollars. The campaign boosted only the visibility of the call for participants. It offered no compensation or incentive of any kind to respondents, and its sole function was to extend the distribution of the recruitment message. No participant received payment, consistent with the informed-consent terms described above.

\section{Data Analysis}

\subsection{Analytic Approach and Its Rationale}

The analytic strategy follows from the nature of the data and the purpose of each research question. It is stated here in full so that every test reported in Chapter~4 can be read against an explicit rationale rather than inferred after the fact. Three principles govern the plan.

First, the survey produces ordinal data. The Likert items record ranked agreement whose intervals are not guaranteed to be equal, and the descriptive results confirm that the composite distributions are non-normal, left-skewed, and concentrated near the top of the scale. Methods that assume normally distributed, interval-scaled data are therefore not the principled default. The analysis accordingly adopts distribution-free procedures as its primary inferential tools: the Mann--Whitney $U$ test for the two-group field comparison, the Kruskal--Wallis test for comparisons across the experience cohorts, Dunn's test with Holm adjustment for the post-hoc contrasts that follow a significant Kruskal--Wallis result, rank-order correlation for the familiarity--value relationship, and the chi-square test of independence for categorical associations. The analysis reports parametric counterparts, principally the independent-samples $t$-test and ordinary least squares regression, alongside the non-parametric results where they aid comparability with the technology-acceptance literature, and never as the sole basis for a conclusion. The choice between a parametric and a non-parametric test is made per comparison, guided by the Shapiro--Wilk test and by visual inspection of the distributions, as the comparative subsections below describe.

Second, because the realized sample is large, statistical significance is easy to obtain and, on its own, uninformative: at this sample size even a negligible difference will reach $p<.05$. The analysis therefore foregrounds effect sizes and confidence intervals, so that the practical magnitude of a difference, not the mere fact of its significance, carries the interpretation. Each comparison reports the effect size matched to its test: the rank-biserial correlation for Mann--Whitney, eta-squared ($\eta^2$) for Kruskal--Wallis, and Cram\'er's $V$ for chi-square, with Cohen's $d$ added for comparability against prior work. The analysis reports proportions with Wilson score intervals, which remain well behaved for the skewed, near-ceiling proportions this instrument produces, where the ordinary normal-approximation interval is unreliable. Where many tests are conducted within a family, the Holm step-down correction controls the family-wise error rate without the excessive conservatism of a simple Bonferroni adjustment. The specific estimators and their formulas are given in the Statistical Formulas subsection below.

Third, and prior to any of the above, inference is only as trustworthy as the data are genuine. A conclusion drawn from automated, coordinated, or non-engaged submissions would be a conclusion about noise rather than about the professional community the study means to characterize. The plan therefore begins by establishing that the analyzed responses are organic, that is, that they were produced by real, engaged professional participants rather than by scripts, coordinated rings, or respondents who did not read the items. The data-preparation procedures that follow, together with the response-authenticity analysis reported in Chapter~4, are the mechanism for that assurance. They do not prove that every retained response is genuine, which no screen can establish, but they are designed to remove the identifiable forms of non-organic responding and to document, rather than assume, what confidence in the remaining data is warranted.

\subsection{Data Preparation and Cleaning}

The analysis plan established a priori that responses missing more than 20\% of the substantive items would be excluded. Before any statistical work, the analysis reviewed survey responses for completeness and data quality, and this review materially shaped the analytic dataset. The survey platform recorded 1{,}600 submissions. A manual screening removed 188 of these as abandoned or aberrant, and a post-hoc response-authenticity check, reported in Chapter~4, removed one further submission as an automated entry, leaving 1{,}411 responses for analysis. Three criteria identified the abandoned or aberrant cases. The first was gross incompleteness: submissions that answered no more than the two Section~1 questions and supplied nothing thereafter. The majority of these recorded no answers at all, and the remainder halted after the opening one or two questions. In either form, such submissions fall well beyond the 20\% missing-data threshold established for exclusion and, more consequentially, provide no data against which the research questions could be evaluated. The second and third criteria concerned implausible completion times. Some submissions were recorded in under three minutes, an interval too brief for considered response to a thirty-item instrument estimated to require approximately ten minutes. Others exceeded five days, a pattern indicating a survey left open in a browser and later abandoned rather than completed in a sitting. The lower cutoff follows response-time screening conventions in web-survey methodology, which flag submissions completed at less than roughly half the median completion time as likely careless responding \parencite{Greszki_Speeding_2015, Leiner_TooFast_2019}; a three-minute completion of an instrument with an estimated ten-minute median amounts to roughly 30\% of that median, comfortably below the conventional threshold. The reasoning is concrete rather than merely conventional: completing thirty items in under three minutes allows fewer than six seconds per item, too little time to read a capability statement, weigh it against one's professional experience, and register a considered judgment. A submission produced at that pace cannot represent the engaged, organic responding the instrument was designed to elicit, whatever its source, and retaining it would substitute mechanical speed for deliberation in the very measures the study reports. The large majority of the pruned responses arrived during the month of May, coinciding with the opening and most heavily promoted phase of recruitment. % Verified: <3-minute cutoff anchored to response-time screening conventions. Leiner (2019, Survey Research Methods 13(3), 229-248, doi:10.18148/srm/2019.v13i3.7403) proposes a relative speed index (median time / respondent time) and flags RSI >= 2.0, i.e., under 50% of median completion time; 3 min on a ~10-min median is ~30% (RSI ~3.3), well past that convention. Greszki, Meyer & Schoen (2015, Public Opinion Quarterly 79(2), 471-503, doi:10.1093/poq/nfu058) analyze removing "too fast" responses defined relative to median response times. The >5-day upper cutoff remains a study-specific abandonment judgment; the speeding literature offers no standard upper-bound convention.

The decision to remove these cases rather than retain them rests on a straightforward methodological judgment. Every composite and comparison in this study depends on substantive responses to the Section~2 through Section~5 items, so a submission that stops after Section~1 contributes nothing to those measures while inflating the apparent sample. Retaining near-empty or time-aberrant records would not have enriched the analysis; it would have introduced noise capable of distorting means, variances, and the planned bias diagnostics, most directly the wave analysis, which keys on the order of submission. Removing 188 such records therefore constitutes a conservative form of data hygiene rather than a selective exclusion of inconvenient cases, because the criteria are content-neutral with respect to how respondents answered the substantive items. After removal, item-level missingness among the retained 1{,}411 respondents is negligible and was handled through available-case analysis for the specific analyses in which a required variable was absent.

\subsection{Descriptive Statistical Analysis}

The analysis calculated descriptive statistics for all survey questions to characterize the distribution of responses. For Likert-scale questions, these included frequency distributions, mean scores, median scores, and standard deviations. For binary and categorical questions, they included frequency counts and percentages within each response category.

\subsubsection{Research Question 1 Analysis}

Analysis addressing Research Question 1 focused on Sections 1 and 2 responses to characterize professional awareness of Digital Engineering capabilities. Key metrics included:

\begin{itemize}
  \item Distribution of familiarity ratings from Question 1.1, with percentage of respondents at each familiarity level
  \item Percentage of respondents reporting professional exposure to Digital Engineering within the past two years (Question 1.2)
  \item Mean agreement scores and standard deviations for understanding of specific capabilities (Questions 2.1-2.6)
  \item Percentage of respondents indicating agreement (score \(\geq\) 4) with each capability statement
\end{itemize}

Chapter~4 reports results with 95\% confidence intervals to indicate the precision of population estimates. For the target sample size of 385 respondents, the margin of error for reported proportions would not exceed 5\%; the realized sample of 1{,}411 responses produced substantially tighter intervals. As noted in the sample size determination, these intervals are conditional on the convenience sample approximating a random draw from the accessible population.

\subsubsection{Research Question 2 Analysis}

Analysis addressing Research Question 2 focused on Sections 3, 4, and 5 responses measuring perceived value of Digital Engineering capabilities. Key metrics included:

\begin{itemize}
  \item Mean agreement scores and response distributions for applicability questions (Questions 3.1-3.6)
  \item Mean agreement scores and response distributions for IT value assessment questions (Questions 4.1-4.4)
  \item Mean agreement scores and response distributions for cybersecurity value assessment questions (Questions 5.1-5.6)
  \item Percentage of respondents indicating organizational investment willingness (Questions 4.5 and 5.7)
\end{itemize}

The analysis interpreted responses indicating agreement (score = 4) or strong agreement (score = 5) with value statements as evidence that professionals perceive potential value in Digital Engineering capabilities.

\subsubsection{Research Question 3 Analysis}

Analysis addressing Research Question 3 focused on questions specifically addressing job performance, compliance requirements, and organizational capability enhancement. Key questions included:

\begin{itemize}
  \item Question 3.5: Digital lifecycle management for compliance and service delivery
  \item Question 3.6: Regulatory or compliance requirements that could benefit from Digital Engineering
  \item Question 4.1: Meaningful value to IT processes
  \item Questions 4.2-4.4: Specific IT operational benefits (cycle time, quality, collaboration)
  \item Question 5.1: Meaningful value to cybersecurity operations
  \item Questions 5.2-5.6: Specific security operational benefits (vulnerability remediation, security posture, threat modeling, collaboration, compliance)
\end{itemize}

The percentage of respondents indicating agreement with these statements provided direct evidence regarding whether professionals believe Digital Engineering could help them in their work.

\subsection{Comparative Analysis}

Comparative analyses examined whether awareness levels and value perceptions differ significantly between professional subgroups and experience levels, directly assessing Proposition P3 (variation by field and experience). Such differences carry important implications for adoption strategies and professional development initiatives.

\subsubsection{Professional Field Comparison}

Independent samples t-tests and Mann--Whitney U tests compared responses between IT professionals and security professionals (based on Question 6.3), serving the field component of Proposition P3. The observed distributional shape, corroborated by Shapiro--Wilk tests and visual inspection, guided the choice of parametric versus non-parametric test. As reported in Chapter~4, the Shapiro--Wilk tests rejected normality for all composites (an expected outcome at this sample size, where the test is highly powered against trivial departures), and the composites are left-skewed and ceiling-concentrated. Mann--Whitney U therefore served as the primary test.

\subsubsection{Experience Level Comparison}

Kruskal--Wallis tests compared the composites across the four experience level categories (Question 6.4) because the composite distributions are non-normal; these comparisons serve the experience component of Proposition P3. Dunn's test with Holm adjustment identified which specific groups differed.

\subsubsection{Association Analysis}

Chi-square tests of independence examined associations between categorical variables, such as the relationship between professional field and prior Digital Engineering exposure (Question 1.2), or between experience level and organizational investment willingness (Questions 4.5, 5.7). These categorical associations, reported with Cram\'er's V, complete the assessment of Proposition P3.

\subsubsection{Familiarity and Value Association}

Rank-order correlation between the familiarity item (Question 1.1) and the value composites, supplemented by ordinary least squares regression of the value composites on familiarity, professional field, and experience level, assesses Proposition P4 (familiarity and value).

\subsection{Composite Score Analysis}

Composite scores average the Likert responses within thematic groupings, providing summary measures of awareness and perceived value:

\subsubsection{Awareness Composite}

Average of Question 1.1 and Questions 2.1-2.6 (seven items measuring familiarity and understanding of Digital Engineering capabilities)

\subsubsection{IT Value Perception Composite}

Average of Questions 3.1, 3.3-3.5, and 4.1-4.4 (eight items measuring perceived value of Digital Engineering for IT operations)

\subsubsection{Security Value Perception Composite}

Average of Questions 3.2, 3.6, and 5.1-5.6 (eight items measuring perceived value of Digital Engineering for Information Assurance and cybersecurity)

The analysis assigned items to composites a priori on theoretical grounds rather than deriving them empirically. In particular, the Section~3 applicability items were allocated between the IT Value and Security Value composites according to the domain each item addresses: items concerning IT operations, environment modeling, and lifecycle value (Questions 3.1, 3.3--3.5) were placed with the IT items, while the items concerning Information Assurance relevance and compliance exposure (Questions 3.2, 3.6) were placed with the security items. Internal consistency was then used to confirm, not to construct, these groupings. Cronbach's alpha assessed the internal consistency of composite scores, with values above 0.70 considered acceptable for research purposes. Had internal consistency proven insufficient, item analysis would have been conducted to identify potentially problematic items; as Chapter~4 reports, all three composites exceeded the threshold, so the composites were carried forward as specified here. Chapter~4 also reports item-level diagnostics for the Awareness composite, which combines a familiarity self-report with capability-agreement items, together with the reasoning for retaining its familiarity item despite those diagnostics.

\subsection{Statistical Significance and Effect Sizes}

All inferential tests used \(\alpha\) = 0.05. However, given the large sample size realized in this study, the analysis reports effect sizes alongside every significance test to assess the practical significance of observed differences. Statistical significance alone can mislead when sample sizes are large; effect sizes reveal whether differences matter in practice. The rank-biserial correlation serves as the primary effect size for Mann--Whitney comparisons and eta-squared (\(\eta^2\)) for Kruskal--Wallis tests, consistent with the non-parametric primary tests. The analysis reports Cohen's d as a supplementary measure for comparability with the technology-acceptance literature, interpreting it against conventional thresholds of 0.2 (small), 0.5 (medium), and 0.8 (large). For chi-square tests, Cram\'{e}r's V serves as the effect size measure.

\subsection{Statistical Formulas}

For transparency and reproducibility, the estimators and test statistics used in the analysis are stated explicitly below. Each Likert item and composite was summarized by the sample mean and the sample standard deviation,
\begin{equation}
  \label{eqn:mean_sd}
  \bar{x} = \frac{1}{n}\sum_{i=1}^{n} x_i,
  \qquad
  s = \sqrt{\frac{1}{n-1}\sum_{i=1}^{n}\left(x_i-\bar{x}\right)^2}.
\end{equation}
\addcontentsline{lofm}{formula}{\protect\numberline{\theequation}Sample mean and standard deviation}
The precision of each mean was expressed through a 95\% confidence interval based on Student's $t$ distribution, while agreement proportions $\hat{p}=k/n$ were bounded by the Wilson score interval, which remains well behaved near the response ceiling,
\begin{equation}
  \label{eqn:ci_wilson}
  \bar{x} \pm t_{0.975,\,n-1}\,\frac{s}{\sqrt{n}},
  \qquad
  \frac{\hat{p}+\dfrac{z^2}{2n}\;\pm\; z\sqrt{\dfrac{\hat{p}\left(1-\hat{p}\right)}{n}+\dfrac{z^2}{4n^2}}}{1+\dfrac{z^2}{n}},
\end{equation}
\addcontentsline{lofm}{formula}{\protect\numberline{\theequation}Confidence interval for the mean and Wilson score interval}
with $z=z_{0.975}=1.96$. Internal consistency for each composite of $k$ items was quantified by Cronbach's $\alpha$, in which $\sigma^2_{j}$ denotes the variance of item $j$ and $\sigma^2_{T}$ the variance of the summed score,
\begin{equation}
  \label{eqn:cronbach}
  \alpha = \frac{k}{k-1}\left(1-\frac{\sum_{j=1}^{k}\sigma^2_{j}}{\sigma^2_{T}}\right).
\end{equation}
\addcontentsline{lofm}{formula}{\protect\numberline{\theequation}Cronbach's $\alpha$ (internal consistency)}
Group differences were summarized by Cohen's $d$, computed with the pooled standard deviation $s_p$, and by the rank-biserial correlation $r_{rb}$ derived from the Mann--Whitney statistic $U$,
\begin{equation}
  \label{eqn:cohend_rrb}
  d = \frac{\bar{x}_1-\bar{x}_2}{s_p},
  \quad
  s_p = \sqrt{\frac{(n_1-1)s_1^2+(n_2-1)s_2^2}{n_1+n_2-2}},
  \qquad
  r_{rb} = 1-\frac{2U}{n_1 n_2}.
\end{equation}
\addcontentsline{lofm}{formula}{\protect\numberline{\theequation}Cohen's $d$, pooled $SD$, and rank-biserial correlation}
The Mann--Whitney test itself compares the rank sum $R_1$ of the first group against its expectation, and the Kruskal--Wallis test generalizes the comparison to $g$ groups, where $R_j$ is the sum of the pooled-sample ranks in group $j$, with the $\eta^2$ statistic serving as its effect size,
\begin{equation}
  \label{eqn:u_h_eta}
  U = R_1 - \frac{n_1(n_1+1)}{2},
  \qquad
  H = \frac{12}{N(N+1)}\sum_{j=1}^{g}\frac{R_j^2}{n_j}-3(N+1),
  \qquad
  \eta^2 = \frac{H-g+1}{N-g}.
\end{equation}
\addcontentsline{lofm}{formula}{\protect\numberline{\theequation}Mann--Whitney $U$, Kruskal--Wallis $H$, and $\eta^2$}
Associations between categorical variables were tested with the chi-square statistic, comparing observed counts $O_{ij}$ against the counts $E_{ij}=R_iC_j/N$ expected under independence, and the magnitude of each association was reported as Cram\'er's $V$,
\begin{equation}
  \label{eqn:chi_cramer}
  \chi^2 = \sum_{i,j}\frac{\left(O_{ij}-E_{ij}\right)^2}{E_{ij}},
  \qquad
  V = \sqrt{\frac{\chi^2}{N\left(\min(r,c)-1\right)}}.
\end{equation}
\addcontentsline{lofm}{formula}{\protect\numberline{\theequation}Chi-square statistic and Cram\'er's $V$}
Multiplicity within each family of related comparisons was controlled by the Holm--Bonferroni step-down procedure, which orders the $m$ unadjusted $p$-values ascending and rejects the hypothesis associated with $p_{(i)}$ while
\begin{equation}
  \label{eqn:holm}
  p_{(i)} \le \frac{\alpha}{m-i+1}.
\end{equation}
\addcontentsline{lofm}{formula}{\protect\numberline{\theequation}Holm--Bonferroni step-down threshold}
The procedure stops at the first $i$ for which the inequality fails and retains all remaining hypotheses. These expressions were implemented in a fully scripted analytical pipeline, so that every statistic reported in Chapter~4 is traceable to its definition and reproducible from the response data.

\subsection{Management of Type I and Type II Errors}

Multiple provisions within the analytical design manage the risk of both Type I errors (false positives, in which the analysis concludes that a difference or association exists when it does not) and Type II errors (false negatives, in which the analysis fails to detect a genuine difference or association). These provisions operate at the levels of sample design, instrument design, and analytical procedure.

\subsubsection{Type I Error Management}

The primary Type I error concern in this study arises from conducting multiple statistical tests across the twenty-three Likert-scaled survey items (organized into the four comparison families used in the analysis: awareness, applicability, IT value, and security value) and multiple subgroup comparisons. When many tests are performed at the \(\alpha\) = 0.05 significance level, the family-wise error rate increases, elevating the probability that at least one test produces a spurious significant result. To manage this inflation, the analytical approach uses the Holm--Bonferroni sequential correction procedure for families of related comparisons. The Holm--Bonferroni method provides stronger Type I error control than unadjusted testing while maintaining greater statistical power than the classical Bonferroni correction, which can be overly conservative when applied to large families of tests.

Comparisons are organized into logical families that follow the structure of the research questions. Awareness items, drawn from Sections 1 and 2, address Research Question 1. The remaining three families all bridge Research Questions 2 and 3: applicability items from Section 3, IT value items from Section 4, and IA value items from Section 5. Within each family, the Holm--Bonferroni correction adjusts significance thresholds to maintain the family-wise error rate at \(\alpha\) = 0.05. Cross-family comparisons, such as the relationship between overall awareness and overall perceived value, are treated as planned comparisons evaluated at the unadjusted \(\alpha\) = 0.05 level, as these represent distinct theoretical questions rather than multiple tests of similar hypotheses.

The emphasis on effect sizes alongside significance testing provides an additional safeguard against Type I error interpretation. Even when statistical tests yield significant p-values after correction, effect size assessment indicates whether observed differences reflect practically meaningful magnitudes rather than trivial differences detected through large sample power.

\subsubsection{Type II Error Management}

The design manages Type II error risk primarily through the chosen sample size. The target sample of 385--450 completed responses provides sufficient statistical power to detect meaningful differences in the primary analyses. For the primary descriptive analyses addressing Research Question 1, the sample size supports margins of error within five percentage points at the 95\% confidence level for reported proportions and means, providing adequate precision to characterize awareness levels and value perceptions.

For comparative analyses between professional subgroups, power considerations informed the oversampling target. The design assumed an approximately even split between IT and IA respondents yielding subgroups of roughly 190--225 each, at which independent samples t-tests achieve power exceeding 0.80 to detect medium effect sizes (Cohen's d = 0.5) at \(\alpha\) = 0.05, and power of approximately 0.83--0.89 for small effects (Cohen's d = 0.3). Effects smaller than d \(\approx\) 0.3 could nonetheless go undetected at the planned subgroup sizes. The realized sample eased this concern: the two focal subgroups comprised 501 IT and 501 IA respondents, far exceeding the planning assumption, so the achieved power to detect even small effects is high. Non-significant subgroup results are nonetheless interpreted as insufficient evidence rather than as evidence of no difference.

For experience level comparisons across four categories, the distribution of respondents across categories may result in unequal cell sizes that reduce power for specific pairwise comparisons. The analytical plan addresses this by reporting confidence intervals for all group means, enabling assessment of overlap and practical significance even when formal tests fail to reach significance. Non-parametric alternatives (Kruskal--Wallis with Dunn's post-hoc test) are used when distributional assumptions are not met, as these tests maintain appropriate Type II error rates under non-normal conditions.

\subsubsection{Instrument-Level Error Reduction}

Survey instrument design contributes to error management through several structural features. Established five-point Likert scales with validated anchor sets reduce measurement error that could contribute to both Type I and Type II errors, so that response variability more closely reflects genuine differences in perception than response format artifacts. Consistent scale formatting across Sections 2 through 5 reduces within-respondent variability attributable to format switching, improving signal-to-noise ratios in comparative analyses.

The inclusion of both broad constructs (overall awareness, overall value perception) and specific sub-constructs (awareness of individual Digital Engineering pillars, value for specific operational domains) enables triangulation of findings. Convergent results across related items strengthen confidence in significant findings (reducing false positive concern), while divergent results prompt examination of whether non-significant findings reflect genuine null effects or insufficient measurement sensitivity (informing Type II error assessment).

Composite score analysis, with internal consistency assessed through Cronbach's alpha, further mitigates measurement error by aggregating multiple indicators of each construct. Composite scores based on reliable item sets exhibit less measurement error than individual items, improving the precision of group comparisons and reducing Type II error risk for analyses involving aggregated constructs.

\section{Reliability and Validity Considerations}

\subsection{Content Validity}

Systematically mapping survey questions to research questions, together with aligning to established frameworks in the Digital Engineering and technology acceptance literature, supports content validity. The survey instrument structure directly addresses the four pillars of Digital Engineering (MBSE, the Digital Thread, Digital Twin, and PLM) as established by INCOSE and adopted by NASA, DoW, and the Intelligence Community. The dissertation committee reviewed the instrument questions during development to assess how well they represent the intended constructs and whether they use appropriate professional terminology familiar to IT and IA practitioners. % Confirmed by author 2026-07-02: content-validity review performed by the dissertation committee.

\subsection{Construct Validity}

Established question formats and scale anchors, drawn from validated instruments in technology acceptance research, support construct validity. The agreement scale format follows conventions from TAM and UTAUT research and has demonstrated validity for measuring technology perceptions across diverse populations. A tripartite structure examining awareness, applicability, and perceived value follows established patterns in technology adoption research that have proven effective across multiple domains.

\subsection{Reliability}

Cronbach's alpha assessed internal consistency reliability for the composite scores. Additionally, the standardized question format and scale anchors across sections support response consistency by presenting respondents with familiar response frameworks throughout the survey. The use of consistent Likert scale response options across Sections 2-5 reduces potential confusion and supports reliable responding.

\subsection{Pilot Testing and Instrument Refinement}

Before formal dissertation proposal development, the survey instrument underwent pilot testing during the Spring 2024 semester. The pilot proceeded as instrument-development work under a formal Institutional Review Board exempt determination, separate from and preceding the exempt determination obtained for the dissertation survey itself; no pilot data are included in the analytic sample reported in Chapter~4. % Confirmed by author 2026-07-02: pilot held a formal IRB exempt determination.
The pilot administered an earlier version of the instrument to IT and IA professionals, serving multiple purposes: evaluating question clarity and comprehension among respondents representative of the target population, assessing completion time and respondent fatigue, identifying ambiguous or confusing question formulations, and generating preliminary data to inform the structural design and thematic organization of the final instrument.

Data collected during the pilot study provided empirical evidence that directly shaped the design and structure of the survey questions presented in this methodology. Pilot results informed several consequential design decisions reflected in the current instrument. The pilot findings drove refinement of question wording to reduce ambiguity in how Digital Engineering concepts were described. The revised wording had to provide sufficient context for respondents unfamiliar with Digital Engineering terminology while avoiding language that could prime respondents toward particular response patterns. Pilot completion patterns motivated adjusting section ordering to place awareness questions before value assessment questions, establishing a cognitive progression from recognition through evaluation. The pilot also validated the estimated completion time and confirmed that the five-point Likert scale format produced adequate response distribution across scale points in the pilot sample, without floor or ceiling effects at that stage. The realized sample proved more concentrated toward agreement than the pilot anticipated, a pattern Chapter~4 reports and treats as a threat to the comparative effect sizes.

The pilot study's contribution extends beyond instrument refinement to substantive validation of the research premise. Pilot data confirmed that the target population includes substantial variation in Digital Engineering awareness levels, validating the foundational assumption that awareness and perceived value remain open empirical questions requiring investigation. The observed variance in pilot responses across both awareness and value dimensions confirmed that the instrument captures meaningful differences among respondents rather than producing uniform response distributions that would limit analytical value.

The pilot testing process complemented the content validity established through systematic mapping of survey questions to research questions and alignment with established Digital Engineering and technology acceptance frameworks. Where the content validity process supported measurement of the intended constructs, pilot testing confirmed that respondents interpret questions as intended and can provide meaningful responses within the designed format. The combination of theoretical grounding through framework alignment and empirical refinement through pilot testing strengthens confidence in the instrument's ability to generate valid, reliable data addressing the research questions.

\subsection{Limitations}

Several limitations should be considered when interpreting study results:

\begin{itemize}
  \item The non-probability sampling approach limits generalizability to the broader population of IT and IA professionals. Results should be interpreted as indicative of perceptions within the accessible population rather than definitive measures of the entire professional community.

  \item Self-selection bias may result in over-representation of professionals with existing awareness or interest in Digital Engineering. Individuals who recognize the term or have prior exposure may be more likely to complete the survey, potentially inflating reported awareness levels.

  \item Social desirability bias may influence responses, particularly regarding perceived value questions. Respondents may indicate greater perceived value than they genuinely hold if they believe Digital Engineering represents a progressive or professionally desirable position.

  \item Self-reported awareness and perceptions may not accurately reflect actual knowledge or organizational capabilities. Respondents may overestimate their familiarity with Digital Engineering concepts or underestimate existing capabilities within their organizations that align with Digital Engineering principles but use different terminology.

  \item The cross-sectional design captures perceptions at a single point in time during a period of rapid evolution in both Digital Engineering practices and enterprise IT methodologies. Results represent a temporal snapshot that may not reflect awareness levels or perceptions six months or a year following data collection, particularly given increasing industry attention to digital transformation initiatives.

  \item Non-target respondents who complete the survey may affect results, though demographic questions enable identification and potential exclusion of responses from outside the target population.

  \item Several question stems bundle related capabilities within a single statement (for example, simulating security scenarios and testing defensive measures), a double-barreled wording pattern that means item-level agreement cannot be attributed to a single clause of the stem.

  \item The survey measures perceived value and anticipated benefits rather than actual experienced benefits. Positive perceptions do not guarantee that Digital Engineering would deliver value if implemented, nor do they indicate organizational readiness for adoption.
\end{itemize}

\subsection{Response Bias Mitigation}

Several design decisions embedded within the survey instrument and recruitment strategy proactively mitigate identified sources of response bias. These measures complement the acknowledged limitations by reducing their anticipated impact on data quality and interpretive validity.

Recruitment messaging mitigates self-selection bias, whereby respondents with prior Digital Engineering awareness disproportionately choose to participate, by explicitly framing the survey as targeting all IT and IA professionals regardless of Digital Engineering familiarity. Recruitment materials emphasize that the research seeks perspectives from professionals across the awareness spectrum and that no prior knowledge of Digital Engineering is required or expected. This framing encourages participation from professionals who might otherwise self-exclude on encountering unfamiliar terminology.

Survey questions in Sections 2 through 5 incorporate minimal contextual descriptions of Digital Engineering capabilities within the question stems, reflecting a deliberate balance between two competing methodological concerns. On one hand, questions devoid of contextual information would measure only pre-existing awareness, excluding respondents unfamiliar with Digital Engineering terminology from providing meaningful responses about perceived value. On the other hand, extensive descriptions risk priming respondents toward positive assessments. The instrument resolves this tension by providing sufficient context for respondents to form an informed judgment about each capability statement without advocating for or characterizing the capability in evaluative terms. Each question describes what a Digital Engineering capability does in neutral, functional language, enabling respondents to assess whether that function would provide value within their professional context based on their own domain expertise. This approach reduces confusion for respondents unfamiliar with Digital Engineering terminology who would otherwise face questions referencing concepts they have not encountered, while avoiding framing that suggests a preferred response direction.

The anonymous survey design mitigates social desirability bias by removing professional reputation concerns that might otherwise influence responses. The absence of personally identifiable information collection, combined with clear communication of anonymity protections in the informed consent notice, encourages candid responses about knowledge gaps and skepticism toward unfamiliar methods.

Two analytical checks address acquiescence bias, the tendency for respondents to agree with statements regardless of content: the response-pattern screen for uniform agreement across all items (straight-lining) described in the Content Requirements section, and the low-familiarity/high-agreement diagnostic that flags respondents who report little familiarity yet agree strongly with value statements. The neutral midpoint (``Neither agree nor disagree'') complements these checks by reducing forced-choice artifacts, enabling genuine expression of uncertainty or indifference rather than forcing agreement or disagreement. The binary and ternary response options for investment willingness questions (Questions 4.5 and 5.7) include an ``Unsure'' option that captures genuine uncertainty without channeling indecisive respondents toward affirmative responses. Because all value items are positively keyed, these measures bound rather than eliminate acquiescence, and the residual risk is acknowledged in the limitations.

The analysis compared early and late respondents (wave analysis) to assess whether response patterns differ systematically between initial and subsequent response waves. Significant differences between early and late respondents may indicate non-response bias, as late respondents are often considered more representative of non-respondents, a rationale Chapter~4 revisits in light of the recruitment history. As reported there, this diagnostic returned a null result, with negligible differences between early and late respondents on all three composites. The recruitment-frame concern developed in Chapter~5 therefore rests on the sampling logic itself rather than on a wave signal. The analysis also compared the demographic composition of the sample against known characteristics of the IT and IA professional population to assess representativeness, with identified disparities reported as potential limitations on generalizability.

An additional analytical safeguard enables post-hoc assessment of potential priming effects from the contextual descriptions embedded in survey questions. Comparing value perception scores (Sections 4 and 5) between respondents reporting high Digital Engineering familiarity (Question 1.1 scores of four or five) and those reporting low familiarity (scores of one or two) provides a diagnostic for priming influence. If both familiarity groups report similarly elevated perceived value, the priming hypothesis gains support, suggesting that question descriptions may have influenced responses independent of prior knowledge. If high-familiarity respondents report substantially higher perceived value than low-familiarity respondents, the measurements more likely reflect genuine perception differences informed by prior knowledge and professional experience rather than question framing effects. This comparison does not eliminate priming concerns but provides empirical evidence for assessing their magnitude and interpretive implications.

\section{Ethical Considerations}

This research was designed in accordance with ethical principles for human subjects research and submitted to the Dakota State University Institutional Review Board (IRB). The IRB issued an exempt review determination under 45~CFR~46.104(d)(2), the category covering research limited to survey procedures (determination DSUIRB-20260421-13EX, April 21, 2026), and that determination was in place before any participant recruitment or data collection began; the first survey response was recorded on April 23, 2026. The full IRB Human Subjects Application is reproduced in Appendix~D. Key ethical considerations include:

\subsubsection{Anonymity} The survey collects no personally identifiable information, ensuring respondent anonymity. Demographic questions are limited to professional field and experience level, which cannot identify individual respondents.

\subsubsection{Voluntary Participation} Participation is voluntary with no consequences for non-participation. Respondents may exit the survey at any time before submission without consequence and may skip any questions they choose not to answer.

\subsubsection{Minimal Risk} The research presents minimal risk to participants, with potential inconvenience limited to time required for survey completion. Risks are similar to those encountered in normal daily internet activity.

\subsubsection{Informed Consent} Informed consent is obtained through the participation notice presented before survey questions, with submission constituting consent.

\subsubsection{Data Protection} Data were stored securely using password protection, multifactor authentication, and encryption at rest and in transit. Data were used solely for research purposes as described in the participation notice.

\subsubsection{Data Use Restrictions} Survey raw data shall not be sold or used for any purpose other than the dissertation research.

\subsection{Researcher Positionality and Data Availability}

The candidate works professionally in systems engineering and Information Assurance and previously developed the reference model for Digital Engineering integration that motivates this investigation~\parencite{Bonar_Hastings_2024}. This proximity to the subject informed the framing of the research questions and the design of the instrument, and it is disclosed here in the interest of transparency. Several features of the study were chosen to limit its influence on the findings. The survey is anonymous; the value items are framed as capability descriptions to be judged by respondents against their own context rather than as advocacy; the analysis foregrounds effect sizes and the planned bias diagnostics over selective emphasis; and the entire analytical pipeline is scripted, so that the reported results are reproducible rather than dependent on discretionary choices. One residual lies outside the reach of these provisions. Every item is affirmatively keyed, so a respondent inclined to agree meets no item on which agreement would register as disconfirmation, and no mitigation listed here alters that property of the instrument. Chapter~5 specifies the reverse-coded and knowledge-verification items a future administration would need to close it. The de-identified response data and the analysis code that generated every table and figure in this dissertation are available from the author on reasonable request, subject to the conditions of the Institutional Review Board approval.

\section{Chapter Summary}

This chapter has presented the research methodology used to investigate professional awareness and perceptions of Digital Engineering capabilities among IT and IA professionals. The research uses a quantitative survey methodology and follows a systems engineering lifecycle approach that supports rigor and traceability throughout the research process. Pilot testing during the Spring 2024 semester refined the survey instrument, which systematically addresses the three research questions through 30 questions organized into six thematic sections covering awareness, understanding, applicability, and perceived value for both IT and IA domains.

The design specified a target sample of 385--450 completed responses to provide sufficient statistical power and a maximum margin of error of 5\% at the 95\% confidence level; the realized sample of 1{,}411 responses exceeded this target threefold. Data analysis used descriptive statistics to characterize awareness levels and value perceptions, with comparative analyses examining differences across professional subgroups and experience levels. Composite scores summarized overall awareness and value perceptions, with internal consistency assessed through Cronbach's alpha. The analytical plan addresses the ordinal-interval debate surrounding Likert scale data through dual-reporting of parametric and non-parametric summaries, and manages Type I error inflation from multiple comparisons through the Holm--Bonferroni sequential correction procedure applied within logical families of related tests.

The study incorporates proactive response bias mitigation through recruitment messaging designed to encourage participation across the awareness spectrum, survey question design that provides minimal contextual information sufficient for informed response without evaluative priming, and planned wave analysis to assess non-response bias. Complementing the acknowledged limitations of non-probability sampling and self-selection, these provisions strengthen confidence in the validity of collected data.

The methodology establishes a rigorous foundation and generates empirical evidence regarding professional perceptions of Digital Engineering that can inform future research and practice in IA and IT service delivery. The systems engineering approach to research design supports traceability between research questions, survey instruments, and analytical approaches, thereby demonstrating the disciplinary rigor of the research.